\chapter{Shells, solids and Boolean operations}\label{ch:solids}

\NF{} describes solids by their boundary (B-rep): a \apiref{Shell} is a set of trimmed faces
(\cref{ch:trimming}) sewn together along shared edges. A shell that is closed and
consistently oriented bounds a solid; it has a volume, a centroid and an inertia tensor, it
can be combined with other solids by Boolean operations, sectioned into layers
(\cref{ch:slicing}), written to STEP (\cref{ch:step}) and meshed (\cref{ch:meshing}).

This chapter covers
\begin{itemize}
\item sewing faces into shells and the B-rep data structure (faces, edges, vertices, edge uses),
\item closedness and orientation checks,
\item area, volume and mass properties,
\item the solid primitives (box, sphere, cylinder, cone and torus),
\item Boolean union, intersection and difference, point classification,
\item face labels that survive all these operations,
\item bodies with several shells (voids, separate lumps), and
\item the limitations of the Boolean operations.
\end{itemize}

\section{Sewing faces into a shell}\label{sec:sew}

\begin{lstlisting}
shell = nf.Shell.sew(faces, tolerance=1e-6)
\end{lstlisting}

\apiref{Shell.sew} takes a list of \code{TrimmedSurface} objects and finds their shared
edges:

\begin{enumerate}
\item The boundary loops of every face are mapped to 3D. Where possible the exact image is
used (straight edges of planar faces, iso-parameter curves such as the end circles and the
seam of a cylinder); otherwise the image is fitted.
\item Two boundary curves form the same \emph{edge} if each lies on the other within the
sewing \code{tolerance} (checked at 7 samples by closest-point projection), in the same or the
opposite direction. Closed curves may start at different points.
\item Degenerate boundary pieces -- the poles of revolved surfaces such as a sphere, where a
whole side of the parameter rectangle maps to one point -- are not edges.
\item The seam of a closed surface (the line where $u=0$ and $u=1$ of a cylinder meet)
becomes an edge used \emph{twice by the same face}.
\item Curves that overlap only partially (T-junctions, e.g. a short edge lying on a long one)
are split at the end points of the other curves.
\end{enumerate}

The tolerance is a model-space length. The default $10^{-6}$ suits faces built exactly in
\NF{}; for imported or fitted faces whose edges disagree by more, increase it -- but keep it
well below the smallest feature size, otherwise unrelated edges may be merged.

\subsection{The B-rep structure}\label{sec:brep}

\begin{center}
\begin{tabularx}{\linewidth}{@{}lX@{}}
\toprule
Member & Meaning \\
\midrule
\code{faces} & list of \code{TrimmedSurface} (index = face number) \\
\code{vertices} & list of \code{Vec3}: end points of the edges \\
\code{edges} & list of \code{ShellEdge} \\
\code{ShellEdge.curve} & the 3D edge curve (from the first face that uses it) \\
\code{ShellEdge.v0}, \code{v1} & start and end vertex index (equal for closed edges such as
circles) \\
\code{ShellEdge.uses} & list of \code{Coedge}: the uses of the edge by face loops \\
\code{Coedge.face} & face index \\
\code{Coedge.loop} & 0 for the outer loop, $k>0$ for hole $k-1$ \\
\code{Coedge.index} & position of the trim curve in the loop \\
\code{Coedge.reversed} & \code{True} if the loop runs against the direction of the edge curve \\
\code{free_edges()} & indices of edges used by only one loop (holes in the shell) \\
\code{non_manifold_edges()} & indices of edges used more than twice \\
\code{components()} & connected components as separate shells \\
\bottomrule
\end{tabularx}
\end{center}

In a \emph{closed}, manifold shell every edge has exactly two uses. In a consistently
oriented shell the two uses traverse the edge in opposite directions (one of them has
\code{reversed=True}): neighbouring faces then have normals pointing to the same side.

\subsection{Example: a cylinder from three faces}

\example{11_shells.py}

\begin{lstlisting}[style=text]
sewn: 3 faces, 3 edges, 2 vertices, closed=True, consistently oriented=False
oriented: consistently oriented=True, volume 6.2831853072 (pi r^2 h = 6.2831853072)
edge 0: closed, length 6.283185, used by [face 0 loop 0, face 1 loop 0 rev]
edge 1: open, length 2.000000, used by [face 0 loop 0, face 0 loop 0 rev]
edge 2: closed, length 6.283185, used by [face 0 loop 0, face 2 loop 0 rev]
cup: closed=False, free edges [2], non-manifold edges [], area 15.70796327
volume of an open shell: GeometryError: shell: volume requires a closed shell
volume 6.28318531, centroid (0.000000, -0.000000, 1.000000)
inertia about the centroid (unit density):
     3.665191    0.000000    0.000000
     0.000000    3.665191   -0.000000
     0.000000   -0.000000    3.141593
analytic: Ixx = m (3 r^2 + h^2) / 12 = 3.665191, Izz = m r^2 / 2 = 3.141593
principal moments (3.141593, 3.665191, 3.665191)
rotated volume 6.28318531, mirrored volume 6.28318531
\end{lstlisting}

The topology is exactly what one expects of a cylinder: the side face (face 0) shares a
closed circular edge with each lid, and its seam (edge 1, from the bottom vertex to the top
vertex) is used twice by the side face itself, once in each direction. Leaving out the top
lid gives an open ``cup'' with one free edge.

\section{Closedness and orientation}\label{sec:shell-orientation}

\begin{center}
\begin{tabularx}{\linewidth}{@{}lX@{}}
\toprule
Query & Meaning \\
\midrule
\code{is_closed} & watertight: every edge used exactly twice \\
\code{is_consistently_oriented} & every edge used twice is traversed in opposite directions
by its two uses \\
\code{oriented()} & a copy with faces flipped so that the orientation is consistent and
closed components point outwards \\
\bottomrule
\end{tabularx}
\end{center}

\code{Shell.sew} keeps the faces as they are; their normals are whatever the surface
parametrisations give. In the cylinder example the side face (an extruded circle) happens to
point outwards while both lids point up, so the raw shell is not consistently oriented.
\code{oriented()} fixes this in two steps:

\begin{enumerate}
\item Within every connected component, faces are flipped by breadth-first propagation over
shared edges until all neighbours agree. A non-orientable configuration (a Möbius-like
arrangement) raises \code{GeometryError}.
\item Closed components are then oriented by \emph{nesting depth}: a component that lies
inside an even number of other closed components is an outer boundary and gets outward
normals (positive volume); one inside an odd number bounds a void and gets normals pointing
into the void (negative volume). See \cref{sec:voids}.
\end{enumerate}

\begin{tip}
Call \code{.oriented()} on every shell you sew yourself. The solid primitives and the results
of Boolean operations are already oriented.
\end{tip}

\section{Area, volume and mass properties}\label{sec:mass}

\begin{center}
\begin{tabularx}{\linewidth}{@{}lX@{}}
\toprule
Method & Meaning \\
\midrule
\code{area(rel_tol=1e-8)} & total area of all faces (also for open shells) \\
\code{volume(rel_tol=1e-8)} & enclosed volume; requires a closed, consistently oriented
shell; positive for outward normals \\
\code{mass_properties(rel_tol=1e-8)} & \apiref{MassProperties} for unit density \\
\bottomrule
\end{tabularx}
\end{center}

The volume is computed with the divergence theorem,
\[
V = \tfrac13 \sum_{\text{faces}} \iint S\cdot(S_u\times S_v)\,du\,dv ,
\]
where each face integral is evaluated over the trimmed region with the same Green's theorem
integration as the trimmed area (\cref{sec:trim-area}). \code{rel_tol} is the relative
accuracy target of the integration. For a closed shell that is not consistently oriented, or
an open shell, \code{volume()} raises \code{GeometryError}.

\apiref{MassProperties} contains

\begin{center}
\begin{tabularx}{\linewidth}{@{}lX@{}}
\toprule
Field & Meaning \\
\midrule
\code{volume} & volume (= mass for unit density) \\
\code{centroid} & centre of mass \\
\code{inertia} & $3\times3$ NumPy array, inertia tensor about the centroid,
$I_{ij} = \int (|r|^2\delta_{ij} - r_i r_j)\,dV$ with $r$ measured from the centroid
(diagonal: moments of inertia; off-diagonal: products of inertia with the minus sign
included) \\
\code{principal_moments} & eigenvalues of \code{inertia}, ascending \\
\code{principal_axes} & $3\times3$ NumPy array of the corresponding unit axes \\
\bottomrule
\end{tabularx}
\end{center}

For a material of density $\rho$ multiply \code{volume} and \code{inertia} by $\rho$. The
integrals are evaluated relative to the centre of the bounding box and combined with the
parallel-axis theorem, which keeps them accurate for parts far from the origin. Voids are
subtracted automatically. \code{mass_properties()} raises \code{GeometryError} if the volume
is not positive (e.g. a shell oriented inside out).

\subsection{Transforming shells}

\code{shell.transformed(T)} (or \code{T * shell}) maps every face surface with the affine
transform \code{T} (trim loops are unchanged, since they live in parameter space) and sews the
faces again. Orientation-reversing maps -- mirrors and negative scale factors -- would turn the
normals inwards; \code{transformed} flips the faces in that case, so outward normals stay
outward and the volume stays positive (see the last line of the example output).
Transforms are described in \cref{ch:transforms}.

\section{Solid primitives}\label{sec:primitives}

Five functions create closed, outward-oriented shells with exact NURBS faces:

\begin{center}
\begin{tabularx}{\linewidth}{@{}>{\raggedright\arraybackslash}p{0.33\linewidth}
  >{\raggedright\arraybackslash}X >{\raggedright\arraybackslash}p{0.2\linewidth}@{}}
\toprule
Function & Shape & Face labels \\
\midrule
\apiref{make_box}\code{(corner, size)} & axis-aligned box, edge lengths \code{size.x},
\code{size.y}, \code{size.z} $>0$; six planar faces & \code{bottom}, \code{top} ($z$),
\code{front}, \code{back} ($y$), \code{left}, \code{right} ($x$) \\
\apiref{make_sphere}\code{(center, radius)} & one revolved face (semicircle about an axis
parallel to $z$), poles at \code{center} $\pm$ radius $\hat z$ & \code{surface} \\
\apiref{make_cylinder}\code{(base, axis, radius, height)} & extruded circle plus two planar
lids trimmed by exact circles; base disc centred at \code{base} & \code{side},
\code{bottom}, \code{top} \\
\apiref{make_cone}\code{(base, axis, r0, r1, height)} & frustum with radius \code{r0} at the
base and \code{r1} at the top; one radius may be 0 (apex as a pole, no lid there) &
\code{side}, \code{bottom}, \code{top} \\
\apiref{make_torus}\code{(center, axis, major_radius, minor_radius)} & one doubly periodic
revolved face; $0 < r < R$ & \code{surface} \\
\bottomrule
\end{tabularx}
\end{center}

Axes are normalised internally. Primitives in other orientations are made with a
non-axis-aligned \code{axis} argument or by transforming the result; a rotated box, for
example, is

\begin{lstlisting}
box = nf.make_box((0, 0, 0), (2, 1, 1)).transformed(
    nf.Transform.rotation((0, 0, 0), (0, 0, 1), math.radians(30)))
\end{lstlisting}

\section{Boolean operations}\label{sec:booleans}

\begin{lstlisting}
result = nf.boolean(a, b, op, tolerance=1e-6)
\end{lstlisting}

\apiref{boolean} combines two closed shells; \code{op} is a \apiref{BooleanOp}:

\begin{center}
\begin{tabularx}{\linewidth}{@{}lX@{}}
\toprule
Operation & Result \\
\midrule
\code{BooleanOp.Union} & $A\cup B$: the parts of $A$ outside $B$ and of $B$ outside $A$ \\
\code{BooleanOp.Intersection} & $A\cap B$: the parts of $A$ inside $B$ and of $B$ inside $A$ \\
\code{BooleanOp.Difference} & $A\setminus B$: the parts of $A$ outside $B$ and of $B$ inside
$A$ (flipped) \\
\bottomrule
\end{tabularx}
\end{center}

The result is again a closed, oriented shell -- possibly with several components (a
difference that cuts a body in two, a union of disjoint bodies, a body with a void), and
possibly empty (the intersection of disjoint bodies, the difference of identical bodies).
The inputs are oriented outwards first, so they do not have to be oriented by the caller.
\code{tolerance} is the length below which overlaps count as contacts (see
\cref{sec:boolean-algorithm}).

\example{11_booleans.py}

\begin{lstlisting}[style=text]
box     : 6 faces ['bottom', 'top', 'front', 'back', 'left', 'right'], volume 24.0000000000 (exact 24.0000000000)
sphere  : 1 faces ['surface'], volume 4.1887902048 (exact 4.1887902048)
cylinder: 3 faces ['side', 'bottom', 'top'], volume 6.2831853072 (exact 6.2831853072)
cone    : 2 faces ['side', 'bottom'], volume 3.1415926536 (exact 3.1415926536)
frustum : 3 faces ['side', 'bottom', 'top'], volume 3.6651914292 (exact 3.6651914292)
torus   : 1 faces ['surface'], volume 9.8696044011 (exact 9.8696044011)
block - bore: 7 faces, closed=True, volume 30.4292036749 (32 - 2 pi 0.25 = 30.4292036732)
labels after the Boolean: ['bottom', 'top', 'front', 'back', 'left', 'right', 'side']
faces labelled 'side' (the bore wall): [6]
spheres Union       : volume 7.068583464, exact 7.068583471, rel. error 9.5e-10
spheres Intersection: volume 1.308996945, exact 1.308996939, rel. error 4.9e-09
spheres Difference  : volume 2.879793259, exact 2.879793266, rel. error 2.2e-09
point_in_shell(drilled, (0, 0, 1)) = False
point_in_shell(drilled, (1.5, 1.5, 1)) = True
point_in_shell(drilled, (0, 0, 3)) = False
slotted pipe: 7 faces, faces per label {'wall': 1, 'inlet': 1, 'outlet': 1, 'slot': 4}
hollow box: 2 components, volumes +64.000000, -4.188790, total 59.811210 (64 - 4/3 pi = 59.811210)
centre inside the solid? False
\end{lstlisting}

The volumes of the primitives are exact to the integration accuracy. Results of Booleans
between curved faces are accurate to about $10^{-9}$ relative: the cut edges are fitted
intersection curves (\cref{sec:isect-ss}), whose fit tolerance limits the accuracy of the
trimmed faces. For the drilled block the volume error is $1.7\times10^{-9}$ in absolute terms.

\subsection{Point classification}

\apiref{point_in_shell}\code{(shell, point, tolerance=1e-6)} returns \code{True} if the
point lies inside the closed shell. It casts rays from the point and counts their crossings
with the trimmed faces; rays that hit a face boundary or graze a face are discarded, and the
answer is the majority of several valid rays. Points within the tolerance of the boundary may
be classified either way. Voids are respected: the centre of the hollow box in the example is
outside the solid.

\subsection{How Booleans work}\label{sec:boolean-algorithm}

\begin{enumerate}
\item \textbf{Intersect.} Every pair of faces with overlapping bounding boxes is intersected
(surface/surface intersection). Each intersection curve is clipped to the parameter intervals
on which both p-curves lie inside their trimmed faces.
\item \textbf{Split.} Every face is split along its cuts by a \emph{planar arrangement} in its
parameter domain: cuts are clipped to the face, all curves are split at their mutual
intersections, dangling pieces are removed and the faces of the resulting graph are traced.
Cuts may end on loops or on other cuts and may cross each other.
\item \textbf{Classify.} Each piece is classified as inside or outside the other shell by
\code{point_in_shell} at an interior point of the piece.
\item \textbf{Select and sew} the pieces according to the operation; pieces of $B$ in a
difference are flipped. The pieces keep the orientation of their outward-oriented shells, so
voids stay voids.
\end{enumerate}

Several special situations are handled explicitly:

\begin{itemize}
\item \textbf{Coincident faces} (two boxes sharing a face, identical solids) are detected by
sampling before the surface/surface intersection. Each face is cut along the other face's
loops; the common part is classified ``on'' with the same or the opposite normal and kept or
dropped according to the following table. So two boxes sharing a face fuse into one solid,
identical solids give themselves (union, intersection) or nothing (difference).

\begin{center}
\begin{tabular}{@{}llll@{}}
\toprule
piece & union & intersection & difference $A-B$ \\
\midrule
$A$, on, same normal & keep & keep & drop \\
$A$, on, opposite normal & drop & drop & keep \\
$B$, on (either) & drop & drop & drop \\
\bottomrule
\end{tabular}
\end{center}

\item \textbf{Tangential contacts} -- spheres touching at a point, a cylinder touching a box
face along a line -- are not cut; touching faces are classified as a whole. The touching lumps
of a union stay separate components.
\item \textbf{Overlaps thinner than the tolerance} count as contacts: two spheres overlapping
by $10^{-7}$ give two touching lumps and an empty intersection. Thicker slivers are cut and
kept, e.g. circle segments of $10^{-5}$.
\end{itemize}

\subsection{Limitations}\label{sec:boolean-limits}

\begin{warning}[Degenerate configurations]
Booleans are numerically sensitive where the inputs meet in degenerate ways. Known cases:
\begin{itemize}
\item \textbf{Intersection curves through poles and along seams} are handled: a box whose
corner lies exactly at the centre of a sphere (its faces contain the poles, one contains the
seam) or a plane along the seam line of a cylinder give closed results on all platforms. Up to
version 0.9.0 the sphere case was closed with Apple Clang but open with GCC and MSVC (issue
\#63): the in/out test of the intersection curve on the seam depended on the sign of rounding
errors.
\item \textbf{Coincident faces whose common part crosses the seam} of a closed surface are
not handled (the projected loop would jump in the parameter domain).
\end{itemize}
Always check \code{result.is_closed} (and, if in doubt, the volume) after a Boolean whose
inputs share poles, seams or tangent faces exactly. Shifting one input by a small, physically
irrelevant amount is usually the simplest remedy.
\end{warning}

Run time grows with the number of face pairs whose boxes overlap and with the complexity of
the intersection curves. Booleans of primitives take fractions of a second to a few seconds;
chains of many operations on complex parts take correspondingly longer. The face pairs are
independent, but the operation runs single-threaded (parallelisation was measured and not
adopted, see the benchmarks page of the online documentation).

\section{Splitting a solid by a plane}\label{sec:split-solid}

Symmetric parts are often simulated as a half or a quarter, with a symmetry boundary condition
on the cut. \apiref{split} with a shell and a \code{Plane} cuts a closed solid into the part
below the plane (opposite to its normal) and the part above:

\begin{lstlisting}
below, above = nf.split(shell, nf.Plane(point, normal), tolerance=1e-6, cut_label="cut")
\end{lstlisting}

Both parts are closed, consistently oriented shells. The new faces lying in the plane carry the
label \code{cut_label}, so they become their own physical group when meshed
(\cref{ch:meshing}); all other faces keep their labels. If the plane misses the solid, one part
is the whole shell and the other an empty shell. Internally the shell is intersected with, and
reduced by, a box that covers it on one side of the plane -- so the same rules and limitations
as for Booleans apply (\cref{sec:boolean-limits}).

\example{11_split_solid.py}

\begin{lstlisting}[style=text]
part volume 7.49734518, halves 3.74867259 + 3.74867259
right half: closed=True, labels ['back', 'bottom', 'front', 'right', 'side', 'symmetry', 'top']
  symmetry face: area 0.600000, outward normal (-1, -0, -0)
  symmetry face: area 0.600000, outward normal (-1, -0, -0)
cap above: 1.18019135 (pi h^2 (3 - h) / 3 = 1.18019164)
rest below: 3.00859856, sum 4.18878992 (4 pi / 3 = 4.18879020)
\end{lstlisting}

The plane through the bore cuts the plate into two halves of equal volume; each half has two
symmetry faces (on either side of the bore) whose outward normals point against the kept
half. The sphere caps deviate from the exact volume by about $2.5\times10^{-7}$ relative,
because the cut circle on the oblique plane is a fitted intersection curve.

\section{Closing planar openings}\label{sec:cap}

A shell built from surfaces -- an extruded profile, a revolved wall, a pipe between two
cylinders -- is often open at planar ends. \apiref{cap} closes such openings with planar
faces:

\begin{lstlisting}
solid = nf.cap(shell, tolerance=1e-6, label="cap")
\end{lstlisting}

The free edges of the shell are chained into closed loops. Every loop must lie in a plane
within \code{tolerance}, otherwise \code{GeometryError} is raised. Loops in the same plane form
one face: nested loops become holes (the annular end of a pipe), loops inside holes become new
faces. Each new face is a planar rectangle trimmed by the edges of the opening, which are
mapped into the rectangle's parameter space by an affine map -- so the trim curves are exact
(a circular opening gives an exact trim circle). The new faces get the label \code{label}, the
result is sewn and oriented outwards. A closed shell is returned unchanged.

\example{11_cap.py}

\begin{lstlisting}[style=text]
pipe wall: 2 faces, closed=False, free edges 4
capped: 4 faces, closed=True, volume 28.27433388 (pi (R^2 - r^2) h = 28.27433388)
  end face: 1 hole, area 9.42477796 (3 pi = 9.42477796)
  end face: 1 hole, area 9.42477796 (3 pi = 9.42477796)
saddle: GeometryError: cap: an opening is not planar (deviation 0.500000 > tolerance)
\end{lstlisting}

\section{Thickening open shells}\label{sec:thicken}

Thin-walled parts -- lens holders, shells of a housing, membranes -- are conveniently modelled
as a surface and given a wall thickness afterwards. \apiref{thicken} turns an open shell into a
solid:

\begin{lstlisting}
solid = nf.thicken(shell, thickness, symmetric=False, tolerance=1e-6, wall_label="wall")
\end{lstlisting}

Every face is offset along its normal $S_u\times S_v$ by \code{thickness} (a negative value
offsets against the normal; with \code{symmetric=True} the faces are offset by
$\pm$\code{thickness}/2, so the original surface becomes the mid-surface). The offset faces keep
the trim loops of their originals, since offsets keep the parametrisation
(\cref{ch:surfaces}). Every free edge of the shell is closed by a ruled side wall between the
edge and its offset. The result is sewn and oriented outwards. Labels: offset faces get the
original label with the suffix \code{_offset} (unlabelled faces: \code{offset}), side walls get
\code{wall_label}.

\begin{warning}[Smooth shells only]
The faces must join tangent-continuously along their shared edges -- then their offsets meet
again. At a crease (the corner of an open box) the offsets separate, the result is not closed
and \code{GeometryError} is raised. The offset also fails where the thickness reaches a radius
of curvature of a face. Thicken creased parts face by face and combine them with Booleans, or
use \apiref{boolean} with an inner body for hollowing.
\end{warning}

\example{11_thicken.py}

\begin{lstlisting}[style=text]
normal at (-5.00, 0.00, 8.66): (+0.500, -0.000, -0.866) - towards the centre, so the wall grows inwards
thick dome: closed=True, faces ['', 'offset', 'wall'], volume 149.356550 (spherical shell sector 149.356551)
washer: volume 3.76991118 (pi (2^2 - 1^2) 0.4 = 3.76991118), walls 2
open box: GeometryError: thicken: the result is not closed (24 free edges) - the faces must join smoothly so that their offsets meet
\end{lstlisting}

Check the direction of $S_u\times S_v$ first, as the example does: for the revolved arc it
points towards the centre of the sphere, so a positive thickness grows the wall inwards.

\section{Pipes}\label{sec:pipes}

Wires, coolant channels, waveguides and flow domains are pipes: a circular cross-section swept
along a path. \apiref{make_pipe} builds the closed solid:

\begin{lstlisting}
pipe = nf.make_pipe(path, radius, end_radius=-1.0, tolerance=1e-6)   # constant or linear taper
pipe = nf.make_pipe(path, lambda t: r(t), tolerance=1e-6)             # radius as a function
\end{lstlisting}

The radius is a number (constant), varies linearly to \code{end_radius}, or is a Python function
of the normalised path parameter $t\in[0,1]$ ($t=0$ at the path start). The wall is the circle
swept with the rotation-minimising frame (as \apiref{sweep}), so it does not twist; its
deviation from the ideal pipe is at most \code{tolerance}. Open paths get planar end faces
(\apiref{cap}). The faces are labelled \code{wall}, \code{inlet} (path start) and
\code{outlet} (path end) -- ready-made physical groups for a flow simulation. A closed path
gives a ring without end faces. The wall surface alone is available as
\apiref{pipe_surface}\code{(path, radius_function, options)}.

The radius must stay positive and smaller than the radius of curvature of the path, otherwise
the wall intersects itself; the radius function should be smooth, since the sections are
skinned with cubic B-splines.

\example{11_pipe.py}

\begin{lstlisting}[style=text]
elbow: closed=True, faces ['inlet', 'outlet', 'wall']
  inlet: area 0.50265482
  outlet: area 0.28274334
  areas pi r^2: 0.50265482 and 0.28274334
bulge: volume 4.99291982 (pi int r^2 dx = 4.99292007)
ring: 1 face, volume 9.47482029 (2 pi^2 R r^2 = 9.47482023)
\end{lstlisting}

The end faces are exact discs (their area is $\pi r^2$ to all digits shown); the volumes agree
with the exact values to about $5\times10^{-8}$ relative, the deviation allowed by the default
wall tolerance.

\section{Defeaturing and hollowing}\label{sec:defeature}

Small holes force fine meshes without affecting the quantity of interest much -- simulation
models are therefore often \emph{defeatured}. \apiref{remove_holes} fills every hole whose loop
is smaller than \code{max_size} across (bounding-box diagonal of its edges):

\begin{lstlisting}
clean = nf.remove_holes(shell, max_size, tolerance=1e-6)
\end{lstlisting}

The hole loops are dropped from their faces (the faces are filled there), and the faces that
formed the feature -- the wall of a through hole, the wall and bottom of a blind hole -- are
then attached to the rest only along free edges; they are removed repeatedly until the shell
is closed again. Features that are not holes (bosses, slots that cut an outer edge) are not
recognised; \code{GeometryError} is raised if the result does not close.

Booleans leave faces split into pieces: the union of two overlapping boxes has three coplanar
top faces. \apiref{merge_faces} joins adjacent faces that carry the same label and lie on the
same surface -- pieces of one split face, or planar faces in one plane with the same normal --
into one face:

\begin{lstlisting}
merged = nf.merge_faces(shell, tolerance=1e-6)
\end{lstlisting}

The merged face is the common surface (or a rectangle in the common plane) trimmed by the union
of the pieces: boundary pieces shared by two faces cancel, the rest forms the new loops, and
holes are kept. The trim curves stay exact, since only affine maps between planar
parametrisations are involved. Faces with different labels are never merged, so physical
groups survive; groups whose boundary does not form valid loops are left as they are.

\apiref{hollow} turns a smooth closed solid into a thin-walled one:

\begin{lstlisting}
part = nf.hollow(shell, thickness, open_labels=[], tolerance=1e-6)
\end{lstlisting}

Without \code{open_labels} every face is offset inwards by \code{thickness} and the offset shell
becomes an inner void (the result has two components, \cref{sec:voids}). With
\code{open_labels} the faces carrying these labels are removed and the rest is thickened
inwards (\cref{sec:thicken}): a cup open at those faces. As for \code{thicken}, the faces must
join smoothly so that their offsets meet; creased solids (a box) raise \code{GeometryError} --
hollow them with a Boolean difference of the solid and a smaller copy instead.

\example{11_defeature_hollow.py}

\begin{lstlisting}[style=text]
part: 10 faces, volume 20.63456859
defeatured: 7 faces, volume 20.85840735 (24 - pi = 20.85840735)
union of two boxes: 14 faces -> merged 10, top face area 7.00000000, volume 7.00000000
hollow sphere: 2 shells, volume 1.13516215 (4/3 pi (1 - 0.9^3) = 1.13516215)
cup: closed=True, labels ['bottom', 'bottom_offset', 'wall'], volume 0.56758107 (2/3 pi (1 - 0.9^3) = 0.56758107)
\end{lstlisting}

The two small holes (diameters 0.4 and 0.5) are filled, the bore of diameter 2 stays: the
defeatured part has exactly the volume of the plate minus the bore.

\section{Fillets and chamfers of edges}\label{sec:edge-blends}

Sharp edges are where fields concentrate in a simulation, and real parts rarely have them.
\apiref{fillet_edges} rounds edges of a solid, \apiref{chamfer_edges} replaces them by a
straight cut:

\begin{lstlisting}
edges = nf.edges_between(shell, "top", "front")      # edges between two labelled faces
rounded = nf.fillet_edges(shell, edges, radius, label="fillet", tolerance=0.0, sizes=[])
bevelled = nf.chamfer_edges(shell, edges, distance, label="chamfer", tolerance=0.0, sizes=[])
\end{lstlisting}

Edges are given by their indices in \code{shell.edges}; \apiref{edges_between} finds them
from the labels of the two adjacent faces. The result is a new closed shell. The new faces are
labelled \code{"fillet"} or \code{"chamfer"}, all other faces keep their labels.

\paragraph{How it works.} The selected edges are grouped by their two adjacent faces. Each
group, which can be a single edge, a chain of edges, or a closed loop such as the rim of a
cylinder or the edge of a bore, gets one blend surface:

\begin{itemize}
  \item \emph{Fillet}: the rolling-ball surface between the two face surfaces
    (\apiref{fillet_surfaces}), with the ball inside the material at a convex edge and outside
    at a concave edge. A concave edge, such as the base of a boss, therefore \emph{gains}
    material.
  \item \emph{Chamfer}: a ruled surface between the two curves at the geodesic distance
    \code{distance} from the edge on both faces (as \apiref{offset_on_surface}). For planar
    faces the chamfer of a right-angled edge removes the triangle with legs \code{distance}.
\end{itemize}

The surfaces are computed on the face surfaces extended beyond their boundaries, so they run
past the ends of the edge. The two faces are then cut back along the contact curves: their
boundary loops are rerouted, not split, which stays exact even where the fillet touches a
boundary edge tangentially. All curves are computed to a twentieth of the sewing tolerance
(\code{tolerance}, by default that of the shell), all groups are built at once on the input
shell, and the result is sewn once.

\paragraph{Ends and corners.} An open chain ends at a vertex where its two faces meet a third
face. What happens there depends on the other selected edges at that vertex:

\begin{itemize}
  \item \emph{Only this chain ends there}, e.g.\ at the end of a single box edge: the blend
    is trimmed by the exact intersection curve with the third face, and that face's boundary is
    rerouted along the same curve. At a convex edge the corner is cut off; at a concave edge
    the end face grows into the corner.
  \item \emph{A second chain ends there and the edge between the remaining faces is smooth}
    (the faces are tangent), e.g.\ the top rim of a plate with rounded corners, where the
    straight and the cylindrical side walls meet tangentially: the two blends continue each
    other and are cut at their common cross-section.
  \item \emph{A second chain ends there at a sharp edge}, e.g.\ two top edges of a box: the
    blends are \emph{mitred}, i.e.\ trimmed by each other along their exact intersection
    curve. Where one of them reaches past the edge between the other faces, it is trimmed by
    the face there, which grows into the corner.
  \item \emph{Three chains end there}, i.e.\ all three edges of the corner are selected:
    every fillet ends at its cross-section through the ball that touches all three faces, and
    the corner is closed by the spherical triangle of that ball between its three contact
    points. The patch is exact and tangent to the three fillets, as on a die. Three chamfers
    end at the straight lines between the points where their contact curves cross, and the
    corner gets a planar triangle (the faces at the corner must be planar).
\end{itemize}

\paragraph{Variable radius.} \code{sizes} gives the radius (or chamfer distance) at points
on the selected edges as a list of \code{(point, size)} pairs. Along each chain the size passes
through these values, blending between neighbours with the quintic smooth step $s_0 + (s_1 -
s_0)(10t^3 - 15t^4 + 6t^5)$ in the arc-length fraction $t$; it is constant before the first
and after the last point (periodic on a closed loop). Chains without points use the
\code{radius} or \code{distance} argument. The cross-sections stay exact circular arcs of the
local radius, in the planes normal to the edge direction $n_1 \times n_2$. Where three fillets
meet at a corner, their radii at the corner ball must agree; a radius given at the corner
vertex holds up to the cross-section through that ball.

\paragraph{Limitations.} A chain must end at a vertex of exactly three faces and three edges,
convex and concave chains cannot meet at a vertex, and a chain cannot continue tangentially
into an edge that is not selected (select the whole tangent chain). \code{GeometryError} is
raised in these cases and for a radius or distance larger than the faces allow, i.e.\ when a
contact curve leaves a face through another edge; its message names the vertex.

\example{11_edge_blends.py}

Output:

\begin{lstlisting}[style=text]
box edge [4]: fillet r = 0.5: 7 faces, closed=True, volume 23.78539817 (24 - 4 r^2 (1 - pi/4) = 23.78539816)
chamfer d = 0.5: volume 23.50000000 (24 - 4 d^2/2 = 23.5), labels ['back', 'bottom', 'chamfer', 'front', 'left', 'right', 'top']
cylinder rim: 4 faces, fillet face area 7.321856
boss base fillet: volume 42.283185 -> 42.412672 (material added)
all edges rounded: 26 faces (20 fillet), volume 22.23598774 (exact 22.23598776)
two chamfers meeting at a corner (mitre): volume 23.16666667 (24 - 7 d^2/2 + d^3/3 = 23.16666667)
variable radius 0.3 -> 0.6: volume 23.81521952 (24 - (1 - pi/4) int r^2 = 23.81521946)
\end{lstlisting}

The fillet face of the cylinder rim is part of a torus. Its area $2\pi(R - r + 2r/\pi)\,\pi r/2
= 7.321856$ for $R = 2$, $r = 0.4$ follows from Pappus' theorem. With all twelve edges of the
box rounded, each corner is an eighth of the ball instead of the sharp corner of the three
edge fillets, so the volume is $24 - 36\,r^2(1 - \pi/4) + 8\,r^3(2 - 7\pi/12)$.

\section{Face labels}\label{sec:face-labels}

Labels give faces stable names across operations -- for boundary conditions, physical groups
in meshes (\cref{ch:meshing}) or named regions in STEP files (\cref{ch:step}).

\begin{center}
\begin{tabularx}{\linewidth}{@{}lX@{}}
\toprule
Member & Meaning \\
\midrule
\code{shell.face_labels} & list with the label of every face (\code{""} for none) \\
\code{shell.with_face_labels(labels)} & a copy with face $f$ labelled \code{labels[f]} (one
label per face) \\
\code{shell.faces_labeled(label)} & indices of the faces with this label \\
\code{face.label}, \code{face.with_label(text)} & label of a single \code{TrimmedSurface} \\
\bottomrule
\end{tabularx}
\end{center}

Labels survive Booleans (every piece of a split face inherits the label of the face it came
from, and new faces from the other operand bring their own labels), transforms,
\code{oriented()}, \code{flipped()}, splitting and the STEP round trip. In the example, the
pipe is labelled \code{wall}, \code{inlet}, \code{outlet} before the slot is cut; after the
Boolean and a translation, the four slot faces carry the label \code{slot} and the original
faces their original names. Because the primitives label their faces automatically, the bore
wall of the drilled block can be found as \code{faces_labeled("side")}.

\begin{tip}
Label faces \emph{before} the operation that creates the final part. The order of faces in a
Boolean result is not specified, but the labels are.
\end{tip}

\section{Bodies with several shells}\label{sec:voids}

A \code{Shell} may consist of several connected components. \code{components()} returns them
as separate shells. Typical cases:

\begin{itemize}
\item \textbf{Voids.} A box minus a sphere lying completely inside it has two components:
the outer box with outward normals (positive volume) and the sphere with normals pointing
into the void (negative volume). The total \code{volume()} is the solid volume,
$64 - \tfrac43\pi$ in the example.
\item \textbf{Separate lumps.} A difference that cuts a body in two, or the union of two
bodies that do not touch, gives several outward-oriented components.
\item \textbf{Touching lumps.} Bodies touching at a point or along a line stay separate
components of the union.
\end{itemize}

The orientation is decided by nesting depth (\cref{sec:shell-orientation}), so hollow
bodies keep their voids through \code{oriented()} and \code{transformed()}, Booleans accept
hollow inputs, \code{mass_properties()} subtracts voids, and the STEP writer exports them as
\code{BREP_WITH_VOIDS}.

\section{Checklist}\label{sec:solids-checklist}

\begin{itemize}
\item After sewing: \code{is_closed}, \code{free_edges()}, \code{non_manifold_edges()}; then
\code{oriented()}.
\item Before measuring: closed and consistently oriented, otherwise \code{volume()} raises.
\item After a Boolean: \code{is_closed}, plausibility of \code{volume()}; beware degenerate
contacts at poles and seams.
\item Use labels to find faces again; do not rely on face indices.
\end{itemize}
